\documentclass[10pt,a4paper]{article}
\usepackage[utf8]{inputenc}
\usepackage[T1]{fontenc}
\usepackage{lmodern}
\usepackage{geometry}
\usepackage{graphicx}
\usepackage{float}
\usepackage{amsmath}
\usepackage{array}
\usepackage{booktabs}
\usepackage[hidelinks]{hyperref}
\usepackage{caption}
\usepackage{enumitem}
\usepackage{titlesec}

\geometry{margin=1.7cm}
\setlength{\parindent}{0pt}
\setlength{\parskip}{3pt}
\setlength{\abovedisplayskip}{6pt}
\setlength{\belowdisplayskip}{6pt}
\setlength{\textfloatsep}{8pt}
\setlength{\floatsep}{6pt}
\setlength{\intextsep}{6pt}
\setlist{itemsep=2pt,topsep=2pt,parsep=0pt,partopsep=0pt}
\captionsetup{font=small,skip=4pt}
\renewcommand{\arraystretch}{0.95}
\titlespacing*{\section}{0pt}{8pt}{4pt}
\titlespacing*{\subsection}{0pt}{6pt}{3pt}
\titlespacing*{\subsubsection}{0pt}{5pt}{2pt}
\linespread{0.96}

\title{NLP Project 2: Language Modelling, Sentiment Classification, and Sentence Boundary Detection}
\author{Kamal Aghazada \and Fatulla Bashirov}
\date{February 2026}

\begin{document}

\maketitle

\begin{abstract}
This report documents Project 2 results for Task 1, Task 2, Task 3, Task 4, and the Extra Task (UI). Tasks 1, 2, and 4 use the shared poetry dataset (\texttt{846} rows, \texttt{9} authors), while Task 3 uses the labeled NLTK \texttt{twitter\_samples} sentiment dataset (\texttt{10,000} tweets). We implemented reproducible pipelines for unsmoothed n-gram modelling, smoothing-method comparison, sentiment classification with multiple algorithms and feature sets, and sentence boundary detection with logistic regression (L1/L2 regularization). The report summarizes methodology, exact metrics, statistical testing, and implementation artifacts.
\end{abstract}

\section{Team Member Contributions}

Work distribution in this project was task-based:
\begin{itemize}
    \item \textbf{Fatulla Bashirov}: implemented and analyzed \textbf{Task 1} (unsmoothed n-gram models and perplexity) and \textbf{Task 2} (smoothing method comparison and ranking).
    \item \textbf{Kamal Aghazada}: implemented and analyzed \textbf{Task 3} (sentiment classification with Naive Bayes, Binary Naive Bayes, Logistic Regression, and significance testing) and \textbf{Task 4} (sentence boundary detection with L1/L2 logistic regression).
\end{itemize}

Both team members reviewed final outputs, validated generated artifacts, and collaborated on report integration and UI delivery.

\section{Introduction and Dataset}

Project 2 uses \texttt{project2/poems\_translated.parquet}, containing the columns \texttt{author}, \texttt{title}, \texttt{url}, \texttt{text}, and \texttt{modern\_text}. All main experiments in this report are based on \texttt{modern\_text}, with deterministic random seed \texttt{42}.

The implemented deliverables are:
\begin{itemize}
    \item Task 1: Unigram, bigram, trigram language models and perplexity.
    \item Task 2: Laplace, Interpolation, Backoff, and Kneser-Ney smoothing comparison.
    \item Task 3: Sentiment classification using Naive Bayes, Binary Naive Bayes, and Logistic Regression with Bag-of-Words and sentiment-lexicon features.
    \item Task 4: Dot-level sentence boundary detection with logistic regression (L1 vs L2).
    \item Extra Task: FastAPI web interface for interactive analysis and inference.
\end{itemize}

\section{Task 1: Unsmoothed N-gram Models and Perplexity}

\subsection{Method}
The Task 1 pipeline (\texttt{project2/task1/task1\_ngram.py}) applies:
\begin{itemize}
    \item Unicode word tokenization with regex \texttt{\textbackslash b\textbackslash w+\textbackslash b} and lowercasing.
    \item Cleaning: remove empty texts and rows without tokens.
    \item Stratified train/test split by author (\texttt{test\_size=0.2}).
    \item Unknown-word policy: tokens with train frequency $<2$ are mapped to \texttt{<unk>}.
    \item Boundary handling for higher-order models with \texttt{<s>} and \texttt{</s>}.
\end{itemize}

Perplexity is computed as:
\[
\mathrm{PPL} = \exp\left(-\frac{1}{N}\sum_{i=1}^{N}\log p(w_i \mid h_i)\right),
\]
and is reported as \texttt{inf} if any zero-probability event occurs in evaluation.

\subsection{Data and Vocabulary Statistics}
\begin{table}[H]
\centering
\caption{Task 1 corpus and split statistics}
\label{tab:task1-stats}
\begin{tabular}{lr}
\toprule
\textbf{Metric} & \textbf{Value} \\
\midrule
Input rows & 846 \\
Rows after cleaning/tokenization & 844 \\
Train rows & 675 \\
Test rows & 169 \\
Train tokens after \texttt{<unk>} mapping & 112,997 \\
Test tokens after \texttt{<unk>} mapping & 29,942 \\
Raw train vocabulary size & 25,098 \\
Retained vocabulary size (freq $\geq 2$) & 10,009 \\
Effective vocabulary with \texttt{<unk>} & 10,010 \\
\bottomrule
\end{tabular}
\end{table}

\subsection{Perplexity Results}
\begin{table}[H]
\centering
\caption{Task 1 model perplexities (unsmoothed MLE)}
\label{tab:task1-ppl}
\begin{tabular}{lrrrr}
\toprule
\textbf{Model} & \textbf{Train PPL} & \textbf{Test PPL} & \textbf{Zero-Prob Test} & \textbf{Unseen Rate (Test)} \\
\midrule
Unigram & 1309.13 & 749.86 & 0 & 0.0000 \\
Bigram & 33.10 & inf & 14,927 & 0.4957 \\
Trigram & 2.33 & inf & 26,193 & 0.8699 \\
\bottomrule
\end{tabular}
\end{table}

Main observation: unsmoothed higher-order models fit training data strongly but fail on held-out text because many n-grams are unseen. This motivates smoothing in Task 2.

Top-frequency outputs (\texttt{project2/task1/results/task1\_top\_ngrams.csv}) also show heavy \texttt{<unk>} presence, confirming the long-tail lexical distribution.

\section{Task 2: Smoothing Method Comparison}

\subsection{Method}
Task 2 (\texttt{project2/task2/task2\_smoothing.py}, \texttt{smoothing\_core.py}) compares:
\begin{itemize}
    \item Laplace (add-one),
    \item Interpolation (bigram and trigram mixtures),
    \item Katz-style Backoff (discount grid),
    \item Kneser-Ney (absolute discounting with continuation probability).
\end{itemize}

Evaluation protocol:
\begin{itemize}
    \item Outer split: train/test with \texttt{test\_size=0.2}.
    \item Inner split: inner-train/validation from outer-train with \texttt{val\_size=0.1}.
    \item Tune hyperparameters on validation perplexity, retrain on final train, evaluate on test.
\end{itemize}

\subsection{Split and Tuning Setup}
\begin{table}[H]
\centering
\caption{Task 2 split configuration}
\label{tab:task2-split}
\begin{tabular}{lr}
\toprule
\textbf{Metric} & \textbf{Value} \\
\midrule
Outer train rows & 675 \\
Inner train rows & 606 \\
Validation rows & 69 \\
Test rows & 169 \\
Interpolation bigram grid & 0.5, 0.6, 0.7, 0.8, 0.9 \\
Interpolation trigram step & 0.1 \\
Discount grid & 0.5, 0.75, 1.0 \\
\bottomrule
\end{tabular}
\end{table}

\subsection{Best Models per Method and Order}
\begin{table}[H]
\centering
\caption{Task 2 best tuned configurations}
\label{tab:task2-methods}
\begin{tabular}{llrrr}
\toprule
\textbf{Method} & \textbf{Order} & \textbf{Best Params} & \textbf{Val PPL} & \textbf{Test PPL} \\
\midrule
Laplace & Bigram & $\alpha=1.0$ & 2186.53 & 3192.02 \\
Interpolation & Bigram & $\lambda_1=0.5,\lambda_2=0.5$ & 429.93 & 642.01 \\
Backoff & Bigram & $d=0.75$ & 410.56 & 600.67 \\
Kneser-Ney & Bigram & $d=0.75$ & 398.88 & 586.84 \\
Laplace & Trigram & $\alpha=1.0$ & 5967.74 & 7850.22 \\
Interpolation & Trigram & $\lambda_1=0.8,\lambda_2=0.2,\lambda_3\approx 0.0$ & 399.42 & 580.93 \\
Backoff & Trigram & $d=0.75$ & 421.50 & 634.05 \\
Kneser-Ney & Trigram & $d=0.75$ & 412.40 & 634.32 \\
\bottomrule
\end{tabular}
\end{table}

\subsection{Ranking Results}
\begin{table}[H]
\centering
\caption{Best method by ranking rule}
\label{tab:task2-ranking}
\begin{tabular}{ll}
\toprule
\textbf{Rule} & \textbf{Best Method} \\
\midrule
\texttt{trigram\_test\_ppl} & Interpolation \\
\texttt{bigram\_test\_ppl} & Kneser-Ney \\
\texttt{avg\_bigram\_trigram\_test\_ppl} & Kneser-Ney \\
\bottomrule
\end{tabular}
\end{table}

By the default decision rule used in the implementation (\texttt{trigram\_test\_ppl}), the selected best method is \textbf{Interpolation}.

\section{Task 3: Sentiment Classification (Naive Bayes, Binary NB, Logistic)}

\subsection{Dataset and Setup}
Task 3 (\texttt{project2/task3/task3\_classification.py}, \texttt{classification\_core.py}) uses the labeled NLTK \texttt{twitter\_samples} corpus:
\begin{itemize}
    \item Total samples: \texttt{10,000} tweets.
    \item Labels: \texttt{positive} and \texttt{negative}.
    \item Split: stratified \texttt{80/20} train/test with seed \texttt{42}.
    \item Train/Test sizes: \texttt{8,000} / \texttt{2,000}.
\end{itemize}

Evaluated classifiers:
\begin{itemize}
    \item Multinomial Naive Bayes,
    \item Binary Naive Bayes,
    \item Logistic Regression.
\end{itemize}

Feature sets:
\begin{itemize}
    \item \textbf{BoW}: CountVectorizer word features,
    \item \textbf{Lexicon}: opinion lexicon counts/ratios + VADER polarity features,
    \item \textbf{BoW+Lexicon}: concatenation of both.
\end{itemize}

\subsection{Main Results}
\begin{table}[H]
\centering
\caption{Task 3 ranking by macro F1 (top configurations)}
\label{tab:task3-summary}
\begin{tabular}{llrrr}
\toprule
\textbf{Rank} & \textbf{Classifier + Features} & \textbf{Accuracy} & \textbf{Macro Precision} & \textbf{Macro F1} \\
\midrule
1 & Logistic Regression + BoW+Lexicon & 0.7990 & 0.7992 & 0.7990 \\
2 & Logistic Regression + BoW & 0.7930 & 0.7938 & 0.7929 \\
3 & Naive Bayes + BoW+Lexicon & 0.7870 & 0.7873 & 0.7869 \\
4 & Binary Naive Bayes + BoW+Lexicon & 0.7850 & 0.7851 & 0.7850 \\
\bottomrule
\end{tabular}
\end{table}

Classifier-level best configurations:
\begin{itemize}
    \item Logistic Regression best: \texttt{BoW+Lexicon}, macro F1 = \texttt{0.79897}.
    \item Naive Bayes best: \texttt{BoW+Lexicon}, macro F1 = \texttt{0.78695}.
    \item Binary Naive Bayes best: \texttt{BoW+Lexicon}, macro F1 = \texttt{0.78499}.
\end{itemize}

\subsection{Statistical Significance Testing}
We used pairwise McNemar tests on the same test predictions (\(\alpha=0.05\)). Key comparisons:
\begin{itemize}
    \item Logistic (BoW+Lexicon) vs Naive Bayes (BoW): \(\chi^2=5.0039\), \(p=0.02529\), significant.
    \item Logistic (BoW+Lexicon) vs Naive Bayes (BoW+Lexicon): \(\chi^2=1.9593\), \(p=0.16159\), not significant.
    \item Logistic (BoW+Lexicon) vs Binary NB (BoW+Lexicon): \(\chi^2=2.4966\), \(p=0.11409\), not significant.
\end{itemize}

\subsection{Analysis}
The best overall model on this split is \textbf{Logistic Regression with BoW+Lexicon}.
However, the significance tests show that its advantage over the strongest Naive Bayes variants is modest and not statistically significant in all pairwise comparisons. Therefore, the most accurate conclusion is:
\begin{itemize}
    \item Logistic Regression is the top-performing classifier in point estimates,
    \item but Logistic and the best Naive Bayes variants are competitive under paired significance testing.
\end{itemize}

\section{Task 4: Sentence Boundary Detection (L1 vs L2)}

\subsection{Method}
Task 4 (\texttt{project2/task4/task4\_sentence\_boundary.py}, \texttt{sentence\_boundary\_core.py}) models each dot character (\texttt{.}) as a binary decision: sentence boundary vs not boundary.

Pipeline highlights:
\begin{itemize}
    \item Dot-level sample extraction from each poem.
    \item Heuristic labeling excludes ellipses and decimal numbers.
    \item 18 engineered features (contextual, positional, lexical, abbreviation and punctuation cues).
    \item Split by poem ID (not by dot) into train/val/test to avoid leakage across dots from the same poem.
    \item Logistic regression with \texttt{saga} solver, comparing L1 and L2 penalties.
    \item Grid search over $C \in \{0.001,0.01,0.1,0.5,1.0,5.0,10.0\}$ on validation F1.
\end{itemize}

\subsection{Dataset Characteristics}
\begin{table}[H]
\centering
\caption{Task 4 dot-level dataset statistics}
\label{tab:task4-data}
\begin{tabular}{lr}
\toprule
\textbf{Metric} & \textbf{Value} \\
\midrule
Total non-empty texts & 844 \\
Total dot samples & 9,428 \\
Positive dots (boundary) & 9,191 \\
Negative dots (not boundary) & 237 \\
Positive rate & 0.9749 \\
Train / Val / Test dots & 6,490 / 875 / 2,063 \\
\bottomrule
\end{tabular}
\end{table}

\subsection{L1 vs L2 Results}
\begin{table}[H]
\centering
\caption{Task 4 final comparison (best $C$ per penalty)}
\label{tab:task4-compare}
\begin{tabular}{lrrrrrr}
\toprule
\textbf{Penalty} & \textbf{Best $C$} & \textbf{Train F1} & \textbf{Val F1} & \textbf{Test F1} & \textbf{Test Acc.} & \textbf{Test Recall} \\
\midrule
L1 & 10.0 & 0.998355 & 0.997585 & 0.983094 & 0.968008 & 0.969192 \\
L2 & 0.01 & 0.997810 & 0.997590 & 0.997477 & 0.995153 & 0.998485 \\
\bottomrule
\end{tabular}
\end{table}

Test-time McNemar comparison between L1 and L2:
\[
\chi^2 = 48.7903,\quad p = 2.8484 \times 10^{-12},
\]
which is statistically significant at $\alpha=0.05$. The selected best penalty is \textbf{L2}.

\subsection{Feature Analysis}
Top absolute coefficients show stable cues across both penalties:
\begin{itemize}
    \item \texttt{is\_ellipsis} (strong negative),
    \item \texttt{next\_is\_quote} (negative),
    \item \texttt{prev\_char\_is\_alpha} (positive),
    \item \texttt{next\_char\_is\_space} (negative),
    \item \texttt{word\_len} (positive).
\end{itemize}

This indicates that punctuation context and immediate local token structure are the strongest predictors for dot-boundary decisions.

\section{Extra Task: FastAPI UI Deliverables}

A unified FastAPI interface was implemented in \texttt{fastapi\_ui/} with a single-page HTML/CSS/JS frontend and API endpoints for all tasks.

Main UI/backend files:
\begin{itemize}
    \item \texttt{fastapi\_ui/main.py}
    \item \texttt{fastapi\_ui/templates/index.html}
    \item \texttt{fastapi\_ui/static/app.js}
    \item \texttt{fastapi\_ui/static/styles.css}
\end{itemize}

The UI provides:
\begin{itemize}
    \item left-side task navigation and per-task Input/Output/Implementation panels,
    \item dataset/split parameter controls,
    \item metric dashboards and comparison tables,
    \item custom text inference (next-word prediction for Tasks 1--2, sentiment prediction for Task 3, sentence detection for Task 4),
    \item confusion matrices and significance tables for classification tasks.
\end{itemize}

To improve responsiveness for heavy tasks, the FastAPI app includes in-memory and persistent disk caching (\texttt{fastapi\_ui/.model\_cache}) and a warmup endpoint (\texttt{/api/warmup}) to prebuild Task 2--4 artifacts.

\section{Discussion and Conclusion}

The Project 2 results show a consistent pattern:
\begin{enumerate}
    \item Unsmoothed higher-order n-gram models are not robust on held-out poetic text, despite strong train fit.
    \item Smoothing is essential. Interpolation gives the best trigram test perplexity, while Kneser-Ney is strongest for bigram and average rules.
    \item For sentiment classification on \texttt{twitter\_samples}, Logistic Regression with BoW+Lexicon gives the best overall macro F1, while top Naive Bayes variants remain competitive in paired significance tests.
    \item Dot-level sentence boundary detection works well with carefully designed local features; L2 regularization significantly outperforms L1 on test data in this setup.
    \item The FastAPI interface makes each task reproducible, interactive, and easier to inspect qualitatively.
\end{enumerate}

Future improvements should include stronger handling of class imbalance in Task 4 and richer domain-specific lexical resources for language modelling and downstream punctuation decisions.

\end{document}
